\documentclass[10pt,a4paper]{article}
\usepackage[utf8]{inputenc}
\usepackage[T1]{fontenc}
\usepackage{lmodern}
\usepackage[margin=1in]{geometry}
\usepackage{xcolor}
\usepackage{amsmath,amssymb,amsfonts}
\usepackage{graphicx}
\usepackage{booktabs}
\usepackage{adjustbox}
\usepackage{tcolorbox}
\usepackage{enumitem}
\usepackage{microtype}
\usepackage{array}
\usepackage{fancyhdr}

% Custom Colors matching house style
\definecolor{bitsnavy}{RGB}{33,53,94}
\definecolor{bitsblue}{RGB}{44,90,160}
\definecolor{bitsgreen}{RGB}{46,125,82}
\definecolor{bitsamber}{RGB}{200,136,30}
\definecolor{bitsred}{RGB}{178,58,72}
\definecolor{lightbg}{RGB}{245,247,250}

% Callout Boxes
\tcolorboxenvironment{intuition}{colback=bitsblue!8,colframe=bitsblue,title=\textbf{Everyday Picture \& Intuition},fonttitle=\bfseries,boxrule=0.8pt,arc=3pt}
\tcolorboxenvironment{worked}{colback=bitsgreen!8,colframe=bitsgreen,title=\textbf{Fully Worked Example},fonttitle=\bfseries,boxrule=0.8pt,arc=3pt}
\tcolorboxenvironment{everyday}{colback=bitsnavy!5,colframe=bitsnavy,title=\textbf{Real-World Picture},fonttitle=\bfseries,boxrule=0.8pt,arc=3pt}
\tcolorboxenvironment{watchout}{colback=bitsred!8,colframe=bitsred,title=\textbf{Watch Out! Pitfalls},fonttitle=\bfseries,boxrule=0.8pt,arc=3pt}
\tcolorboxenvironment{keytake}{colback=bitsamber!10,colframe=bitsamber,title=\textbf{Key Takeaways},fonttitle=\bfseries,boxrule=0.8pt,arc=3pt}

\newenvironment{intuition}{\begin{tcolorbox}}{\end{tcolorbox}}
\newenvironment{worked}{\begin{tcolorbox}}{\end{tcolorbox}}
\newenvironment{everyday}{\begin{tcolorbox}}{\end{tcolorbox}}
\newenvironment{watchout}{\begin{tcolorbox}}{\end{tcolorbox}}
\newenvironment{keytake}{\begin{tcolorbox}}{\end{tcolorbox}}

\newenvironment{steps}{\begin{enumerate}[label=\textbf{Step \arabic*:},leftmargin=*]}{\end{enumerate}}

\newcommand{\secsub}[1]{\noindent\textit{\large #1}\vspace{0.8em}}
\newcommand{\housefig}[2]{
\begin{figure}[htbp]
\centering
\includegraphics[width=0.85\linewidth]{#1}
\caption{#2}
\end{figure}
}

\pagestyle{fancy}
\fancyhf{}
\rhead{\color{bitsnavy}\bfseries WILP, BITS Pilani}
\lhead{\color{bitsnavy}\bfseries ISM Session 11A Study Companion}
\rfoot{Page \thepage}

\begin{document}

% Banner
\begin{tcolorbox}[colback=bitsnavy,colframe=bitsnavy,arc=0pt,outer arc=0pt,top=12pt,bottom=12pt]
\color{white}
\centering
{\large\bfseries WILP, BITS Pilani}\\[4pt]
{\LARGE\bfseries Session 11A: Two-Way ANOVA, Correlation \& Regression}\\[4pt]
{\small Course: Introduction to Statistical Methods (ISM) \quad|\quad Instructor: Dr. YVK Ravi Kumar}
\end{tcolorbox}

\vspace{1em}

\section{Introduction to Two-Way ANOVA}
\secsub{Isolate two independent factor effects at once without interference.}

When analyzing experimental data, we often test how two categorical factors affect a continuous response. Two-Way Analysis of Variance (ANOVA) separates total variation into row factor effects, column factor effects, and unexplained random error.

\begin{intuition}
Imagine testing baked cake height. One factor is oven temperature (low, medium, high). The second factor is sugar brand (Brand A, Brand B). Two-Way ANOVA tells us if oven temperature matters, if sugar brand matters, and keeps them from confusing each other.
\end{intuition}

\subsection{Mathematical Framework}
Total variation ($SST$) is partitioned into sum of squares for row factor ($SSA$), column factor ($SSB$), and error ($SSE$):
\begin{equation}
SST = SSA + SSB + SSE
\end{equation}
The degrees of freedom partition accordingly: $rn - 1 = (r - 1) + (c - 1) + (r - 1)(c - 1)$. Mean squares are calculated as $MS = \frac{SS}{df}$, leading to test statistics $F_{\text{row}} = \frac{MSA}{MSE}$ and $F_{\text{col}} = \frac{MSB}{MSE}$.

\housefig{figures/anova_partition.png}{ANOVA divides total dataset variation into row factor, column factor, and residual error components.}

\begin{worked}
Consider a study with 3 rows and 4 columns. Suppose $SSA = 120$, $SSB = 180$, and $SSE = 60$.
\begin{steps}
\item Calculate degrees of freedom: $df_A = 3 - 1 = 2$, $df_B = 4 - 1 = 3$, $df_E = 2 \times 3 = 6$.
\item Calculate Mean Squares:
\begin{align*}
MSA &= \frac{120}{2} = 60 \\
MSB &= \frac{180}{3} = 60 \\
MSE &= \frac{60}{6} = 10
\end{align*}
\item Compute F-ratios: $F_{\text{row}} = \frac{60}{10} = 6.0$, $F_{\text{col}} = \frac{60}{10} = 6.0$.
\end{steps}
\end{worked}

\section{Covariance and Karl Pearson Correlation}
\secsub{Measure the direction and strength of linear relationship between two variables.}

Covariance measures whether two variables move together in the same direction. Because covariance depends on measurement scale, Karl Pearson standardized it into correlation coefficient $r$, which ranges strictly from $-1$ to $+1$.

\begin{intuition}
Covariance is like measuring combined height and weight in inches and pounds; the raw number changes if you switch to centimeters and kilograms. Correlation converts this number to a standard score between $-1$ and $+1$.
\end{intuition}

\subsection{Formula for Karl Pearson's $r$}
\begin{equation}
r = \frac{\sum (X - \bar{X})(Y - \bar{Y})}{\sqrt{\sum (X - \bar{X})^2 \sum (Y - \bar{Y})^2}} = \frac{\text{Cov}(X,Y)}{s_X s_Y}
\end{equation}

\housefig{figures/correlation_scatter.png}{Scatter plot showing strong positive correlation versus strong negative correlation.}

\begin{worked}
Let us evaluate student anxiety ($X$) versus test scores ($Y$) across 5 students:
\begin{center}
\begin{tabular}{cccccc}
\toprule
Student & $X$ & $Y$ & $(X-\bar{X})$ & $(Y-\bar{Y})$ & $(X-\bar{X})(Y-\bar{Y})$ \\
\midrule
1 & 2 & 90 & -3 & +20 & -60 \\
2 & 4 & 80 & -1 & +10 & -10 \\
3 & 5 & 70 & 0 & 0 & 0 \\
4 & 6 & 60 & +1 & -10 & -10 \\
5 & 8 & 50 & +3 & -20 & -60 \\
\bottomrule
\end{tabular}
\end{center}
\begin{steps}
\item Calculate means: $\bar{X} = 5$, $\bar{Y} = 70$.
\item Sum of cross-products: $\sum (X-\bar{X})(Y-\bar{Y}) = -140$.
\item Sum of squared deviations: $\sum (X-\bar{X})^2 = 20$, $\sum (Y-\bar{Y})^2 = 1000$.
\item Compute $r$:
\begin{equation}
r = \frac{-140}{\sqrt{20 \times 1000}} = \frac{-140}{\sqrt{20000}} = \frac{-140}{141.42} = -0.99
\end{equation}
Conclusion: Nearly perfect strong negative correlation. Higher anxiety correlates with lower test scores.
\end{steps}
\end{worked}

\section{Simple Linear Regression Analysis}
\secsub{Predict a continuous target variable using a straight-line trend model.}

Simple linear regression models the target variable $Y$ as a linear function of predictor $X$ plus random noise: $Y = a + bX + \epsilon$.

\begin{intuition}
Think of simple regression as drawing a baseline path through scattered data points so that the total missed vertical distances are as small as possible.
\end{intuition}

\subsection{Least Squares Formulas}
Slope coefficient $b$ and intercept $a$ minimize squared errors:
\begin{align}
b &= \frac{\sum (X - \bar{X})(Y - \bar{Y})}{\sum (X - \bar{X})^2} \\
a &= \bar{Y} - b\bar{X}
\end{align}

\housefig{figures/regression_line.png}{Fitted least-squares regression line estimating weekly movie gross revenue from advertising spend.}

\begin{worked}
A theater owner records advertising expenditure ($X$ in \$1,000s) and weekly gross revenue ($Y$ in \$1,000s):
Sample points $(X,Y)$: $(2, 9), (5, 16), (4, 15), (2.5, 11), (3, 13), (4, 14), (2.5, 10), (3, 12)$.
Total sample size $n = 8$, $\bar{X} = 3.25$, $\bar{Y} = 12.50$.
Given $\sum (X - \bar{X})^2 = 8.50$ and $\sum (X - \bar{X})(Y - \bar{Y}) = 20.00$.
\begin{steps}
\item Calculate slope $b$:
\begin{equation}
b = \frac{20.00}{8.50} = 2.353 \approx 2.36
\end{equation}
\item Calculate intercept $a$:
\begin{equation}
a = 12.50 - (2.353 \times 3.25) = 12.50 - 7.647 = 4.853 \approx 4.85
\end{equation}
\item Fitted regression model: $\hat{Y} = 4.85 + 2.36X$.
\item Forecast revenue for ad spend $X = 4.5$ (\$4,500):
\begin{equation}
\hat{Y} = 4.85 + 2.36(4.5) = 4.85 + 10.62 = 15.47 \quad (\text{or \$15,470})
\end{equation}
\end{steps}
\end{worked}

\begin{watchout}
Correlation does not mean causation! High correlation between two variables does not prove that changing one directly causes changes in the other.
\end{watchout}

\begin{keytake}
\begin{itemize}
\item Two-Way ANOVA isolates row and column effects simultaneously.
\item Karl Pearson's $r$ measures linear association strength on a scale from $-1$ to $+1$.
\item Least-squares regression minimizes squared vertical errors to yield optimal prediction line parameters $a$ and $b$.
\end{itemize}
\end{keytake}

\end{document}
